\documentclass[a4paper,oneside,14pt]{extarticle}

\usepackage{cmap} % Улучшенный поиск русских слов в полученном pdf-файле
\usepackage[T2A]{fontenc} % Поддержка русских букв
\usepackage[utf8]{inputenc} % Кодировка utf8
\usepackage[english,russian]{babel} % Языки: русский, английский

\usepackage[14pt]{extsizes}

\usepackage{graphicx}
\usepackage{multirow}

\usepackage{tikz}
\usetikzlibrary{shapes, shapes.geometric, arrows, arrows.meta, positioning}

\usepackage{caption}
\captionsetup{labelsep=endash}
\captionsetup[figure]{name={Рисунок}}

\usepackage{amsmath}
\usepackage{amsfonts}

\usepackage{geometry}
\geometry{left=30mm}
\geometry{right=10mm}
\geometry{top=20mm}
\geometry{bottom=20mm}

\usepackage{enumitem}

\usepackage{tabularx}
\usepackage{longtable}
\usepackage{adjustbox}
\usepackage{threeparttable}

% Переопределение стандартных \section, \subsection, \subsubsection по ГОСТу;
\usepackage{titlesec}[explicit]
\titleformat{name=\section,numberless}[block]{\normalfont\large\bfseries\centering}{}{0pt}{}
\titleformat{\section}[block]{\normalfont\large\bfseries}{\thesection}{1em}{}
\titlespacing\section{\parindent}{*4}{*4}

\titleformat{\subsection}[hang]
{\bfseries\large}{\thesubsection}{1em}{}
\titlespacing\subsection{\parindent}{*2}{*2}

\titleformat{\subsubsection}[hang]
{\bfseries\large}{\thesubsubsection}{1em}{}
\titlespacing\subsubsection{\parindent}{*2}{*2}

\usepackage{url}

% Переопределение их отступов до и после для 1.5 интервала во всем документе
\usepackage{setspace}
\onehalfspacing % Полуторный интервал
\frenchspacing
\setlength\parindent{1.25cm}

\usepackage{indentfirst} % Красная строка

% Настройки оглавления
\usepackage{xcolor}
\usepackage{multirow}

% Гиперссылки
\usepackage[pdftex]{hyperref}
\hypersetup{hidelinks}

% Дополнительное окружения для подписей
\usepackage{array}
\newenvironment{signstabular}[1][1]{
	\renewcommand*{\arraystretch}{#1}
	\tabular
}{
	\endtabular
}

\usepackage{enumitem} 
\setenumerate[0]{label=\arabic*)} % Изменение вида нумерации списков
\renewcommand{\labelitemi}{---}

% Листинги 
\usepackage{courier}
\usepackage{listings}
\usepackage{chngcntr} % Listings counter within section set main.tex after begin document
\usepackage{float} % Place figures anywhere you want; ignore floating
\floatstyle{plaintop}
\newfloat{code}{H}{myc}

% Для листинга кода:
\lstset{
	basicstyle=\small\ttfamily,
	language=prolog,
	numbers=left,
	numbersep=5pt,
	% stepnumber=1,
	xleftmargin=17pt,
	% showstringspaces=false,
	numbersep=5pt,
	frame=single,
	tabsize=4,
	captionpos=b,
	breaklines=true,
	breakatwhitespace=true,
	escapeinside={\#*}{*)},
	inputencoding=utf8x,
	backgroundcolor=\color{white},
	% numberstyle=,%\tiny
	keywordstyle=\color{blue},
	stringstyle=\color{red!90!black},
	commentstyle=\color{green!50!black}
}
\lstset{
 morekeywords={domains, predicates, clauses, goal}
}

\lstset{
	literate=
	{а}{{\selectfont\char224}}1
	{б}{{\selectfont\char225}}1
	{в}{{\selectfont\char226}}1
	{г}{{\selectfont\char227}}1
	{д}{{\selectfont\char228}}1
	{е}{{\selectfont\char229}}1
	{ё}{{\"e}}1
	{ж}{{\selectfont\char230}}1
	{з}{{\selectfont\char231}}1
	{и}{{\selectfont\char232}}1
	{й}{{\selectfont\char233}}1
	{к}{{\selectfont\char234}}1
	{л}{{\selectfont\char235}}1
	{м}{{\selectfont\char236}}1
	{н}{{\selectfont\char237}}1
	{о}{{\selectfont\char238}}1
	{п}{{\selectfont\char239}}1
	{р}{{\selectfont\char240}}1
	{с}{{\selectfont\char241}}1
	{т}{{\selectfont\char242}}1
	{у}{{\selectfont\char243}}1
	{ф}{{\selectfont\char244}}1
	{х}{{\selectfont\char245}}1
	{ц}{{\selectfont\char246}}1
	{ч}{{\selectfont\char247}}1
	{ш}{{\selectfont\char248}}1
	{щ}{{\selectfont\char249}}1
	{ъ}{{\selectfont\char250}}1
	{ы}{{\selectfont\char251}}1
	{ь}{{\selectfont\char252}}1
	{э}{{\selectfont\char253}}1
	{ю}{{\selectfont\char254}}1
	{я}{{\selectfont\char255}}1
	{А}{{\selectfont\char192}}1
	{Б}{{\selectfont\char193}}1
	{В}{{\selectfont\char194}}1
	{Г}{{\selectfont\char195}}1
	{Д}{{\selectfont\char196}}1
	{Е}{{\selectfont\char197}}1
	{Ё}{{\"E}}1
	{Ж}{{\selectfont\char198}}1
	{З}{{\selectfont\char199}}1
	{И}{{\selectfont\char200}}1
	{Й}{{\selectfont\char201}}1
	{К}{{\selectfont\char202}}1
	{Л}{{\selectfont\char203}}1
	{М}{{\selectfont\char204}}1
	{Н}{{\selectfont\char205}}1
	{О}{{\selectfont\char206}}1
	{П}{{\selectfont\char207}}1
	{Р}{{\selectfont\char208}}1
	{С}{{\selectfont\char209}}1
	{Т}{{\selectfont\char210}}1
	{У}{{\selectfont\char211}}1
	{Ф}{{\selectfont\char212}}1
	{Х}{{\selectfont\char213}}1
	{Ц}{{\selectfont\char214}}1
	{Ч}{{\selectfont\char215}}1
	{Ш}{{\selectfont\char216}}1
	{Щ}{{\selectfont\char217}}1
	{Ъ}{{\selectfont\char218}}1
	{Ы}{{\selectfont\char219}}1
	{Ь}{{\selectfont\char220}}1
	{Э}{{\selectfont\char221}}1
	{Ю}{{\selectfont\char222}}1
	{Я}{{\selectfont\char223}}1
}

% Работа с изображениями и таблицами; переопределение названий по ГОСТу
\usepackage{caption}
\captionsetup[figure]{name={Рисунок},labelsep=endash}
\captionsetup[table]{singlelinecheck=false, labelsep=endash}

\usepackage[justification=centering]{caption} % Настройка подписей float объектов	

\usepackage{csvsimple}

\usepackage{ulem} % Нормальное нижнее подчеркивание
\usepackage{hhline} % Двойная горизонтальная линия в таблицах
\usepackage[figure,table]{totalcount} % Подсчет изображений, таблиц
\usepackage{rotating} % Поворот изображения вместе с названием
\usepackage{lastpage} % Для подсчета числа страниц

\makeatletter
\renewcommand\@biblabel[1]{#1.} % [1] -> 1. in bibliography
\makeatother

\usepackage{ragged2e} % Перенос слов на следующую строку
\usepackage{pdfpages}

\usepackage{blindtext}

% \usepackage[
%     backend=biber,
% 	style=gost-numeric,
% 	% style=numeric-comp,
% 	language=auto,
% 	autolang=other,
% 	sorting=none
% ]{biblatex}
% \addbibresource{bibliography.bib}
% \usepackage{xparse} % \NewDocumentCommand for creating custom commands
% \NewDocumentCommand{\printbib}{m}
% {\printbibliography[title={#1}]\addcontentsline{toc}{section}{#1}}


\begin{document}

\def\coursename{Функциональное и логическое программирование}
\def\labnumber{5}
\def\myname{Рунов К.А.}
\def\mygroup{ИУ7-64Б}
\def\teachers{Толпинская Н. Б., Строганов Ю. В.}

\begin{titlepage}
    \newgeometry{left=2cm, right=1cm, top=2.5cm, bottom=2.5cm}
    \fontsize{12pt}{12pt}\selectfont

    \noindent
    \begin{center}
        % \fbox
        % \begin{tabular}{|l|r|}
        % {
            % \hline
            \begin{minipage}{0.12\textwidth}
                \includegraphics[width=\linewidth]{img/bmstu_logo.jpg}
            \end{minipage}
            \hfill
            % &
            % \hspace{0.2cm}
            \begin{minipage}{0.85\textwidth}\centering\bfseries
                {
                    \linespread{1}\selectfont
                    \vspace{0.1cm}
                    % \textsc
                    {Министерство науки и высшего образования~Российской~Федерации}

                    % \textsc
                    {Федеральное~государственное~бюджетное~образовательное~учреждение высшего образования}

                    % \textsc
                    {<<Московский государственный технический университет имени~Н.~Э.~Баумана (национальный~исследовательский~университет)>>}

                    % \textsc
                    {(МГТУ им. Н.~Э.~Баумана)}
                    \vspace{0.1cm}
                }
            \end{minipage}
            % \\
            % \hline
        % }
        % \end{tabular}

        \vspace{0.2cm}
        \rule{\linewidth}{2.8pt}
        \rule[3ex]{\linewidth}{1pt}

        \begin{flushleft}
            {ФАКУЛЬТЕТ \uline{<<Информатика и системы управления>> \hfill}}

            \vspace{0.5cm}

            {КАФЕДРА \uline{<<Программное обеспечение ЭВМ и информационные технологии>> \hfill}}
        \end{flushleft}

        % \vspace{1cm}
        \vfill

        {
            \Large{\textbf{
                % \bolduline
                {ОТЧЕТ ПО ЛАБОРАТОРНОЙ РАБОТЕ №\labnumber}
            }}

            \Large{\textbf{
                % \bolduline
                {по курсу <<\coursename>>}
            }}

            \vspace{0.5cm}
        }

        \vspace{0.5cm}

        % \setstretch{1.5}
        % \begin{tabular}{p{\textwidth}}
        %     \uline
        %     {
        %         Разработка программы для моделирования столкновений объектов
        %         в виртуальном пространстве. \hfill
        %     }
        %     \rule{\linewidth}{0.4pt}
        % \end{tabular}

        \fontsize{14pt}{14pt}\selectfont

        \vfill

        \begin{flushleft}
            % {Студент \uline{Рунов Константин Алексеевич \hfill}}
            {Студент \uline{\myname \hfill}}

            \vspace{0.5cm}

            {Группа \uline{\mygroup \hfill}}

            \vspace{0.5cm}

            {Оценка (баллы) \uline{\hfill}}

            \vspace{0.5cm}

            {Преподаватели \uline{\teachers \hfill}}

            % \vspace{0.5cm}

            % {Преподаватель \hfill \ulinetext[4cm]{(Подпись, дата)}{} \ulinetext[4cm]{}{Строганов Д.~В.}}

            \vspace{0.5cm}

        \end{flushleft}

        \vfill

        \the\year\ г.

    \end{center}
\end{titlepage}


\setcounter{page}{2}
\renewcommand{\contentsname}{СОДЕРЖАНИЕ}
% \tableofcontents

% -------------------------------------------------------------------

\begin{lstlisting}
; 1. Напишите функцию, которая уменьшает на 10 все числа из списка-аргумента этой функции, проходя по верхнему уровню списковых ячеек. (Список смешанный структурированный)
(defun f1-helper (x)
  (cond ((numberp x) (- x 10))
        (t x)))

(defun f1 (lst)
  (mapcar #'f1-helper lst))
\end{lstlisting}

\begin{lstlisting}
; 2. Написать фукнцию которая получает как аргумент список чисел, а возвращает список квадратов этих чисел в том же порядке
(defun f2-helper (x)
  (cond ((numberp x) (* x x))
        (t x)))

(defun f2 (lst)
  (mapcar #'f2-helper lst))
\end{lstlisting}

\begin{lstlisting}
; 3. Напишите фукнцию, которая умножает на заданное число-аргумент все числа из заданного списка-аргумента, когда
; а) все элементы списка -- числа,
; б) элементы списка -- любый объекты.
(defun f3-helper (x c)
  (cond ((numberp x) (* x c))
        (t x)))

(defun f3 (lst c)
  (mapcar #'(lambda (x) (f3-helper x c)) lst))
\end{lstlisting}

\begin{lstlisting}
; 4. Написать функцию, которая по своему списку-аргументу lst определяет является ли он палиндромом (то есть равны ли lst и (reverse lst)), для одноуровневого смешанного списка
(defun f4 (lst)
  (equal lst (reverse lst)))
\end{lstlisting}

\begin{lstlisting}
; 5. Используя функционалы, написать предикат set-equal, который возвращает t, если два его множества-аргумента (одноуровневые списки) содержат одни и те же элементы, порядок которых не имеет значения
(defun set-equal-helper (lst1 lst2 flag-call-first-time)
  (cond ((and (null lst1) flag-call-first-time) nil)
        ((null lst1) t)
        ((member (car lst1) lst2)
         (set-equal-helper (cdr lst1) lst2 nil))
        (t nil)))

(defun set-equal (lst1 lst2)
  (and (set-equal-helper lst1 lst2 t)
       (set-equal-helper lst2 lst1 t)))
\end{lstlisting}

\begin{lstlisting}
; 6. Напишите функицю, select-between, которая из списка-аргумента, содержащего только числа, выбирает только те, которые расположены между двумя указанными числами - границами-аргументами и возвращает их в виде списка (упорядоченного по возрастанию (+2 балла))
(defun push-ordered (lst elem)
  (cond ((null lst) (list elem))
        ((< elem (car lst)) (cons elem lst))
        (t (cons (car lst) (push-ordered (cdr lst) elem)))))

(defun select-between-helper (lst start end result-lst)
  (cond ((null lst) result-lst)
        ((and (< start (car lst)) (> end (car lst)))
         (select-between-helper (cdr lst) start end
                            (push-ordered result-lst (car lst))))
        (t (select-between-helper (cdr lst) start end result-lst))))

(defun select-between (lst start end)
  (select-between-helper lst start end nil))
\end{lstlisting}

\begin{lstlisting}
; 7. Написать функцию, вычисляющую декартово произведение двух своих списков-аргументов (напомним, что A x B это множество всевозможных пар (a b), где a принадлежит A, b принадлежит B)
(defun decart (lstA lstB)
  (mapcan #'(lambda (a)
              (mapcar #'(lambda (b)
                          (list a b)) lstB)) lstA))
\end{lstlisting}

\begin{lstlisting}
; 8. Почему так реализовано reduce, в чем причина?
; (reduce #'+ ()) -> 0
; (reduce #'* ()) -> 1
reduce принимает на вход функцию, принимающую либо 0, либо 2 аргумента.
Если reduce передана бинарная функция fn, то
(reduce #'fn '(arg1 arg2 arg3 ... argN)) <=>
<=> (fn (... (fn (fn arg1 arg2) arg3) ...) argN)
Пример:
(reduce #'list '(1 2 3 4)) ; (((1 2) 3) 4)

Если reduce передана функция fn, не принимающая аргументов, то
(reduce #'fn '()) <=> (fn)
Пример:
(reduce #'list '()) ; (list) = NIL
(reduce #'+ '()) ; (+) = 0
(reduce #'* '()) ; (*) = 1

Такое поведение может быть полезно, когда нужно гарантировать, что в случае передачи пустой последовательности, будет возвращено значение по умолчанию, определённое логикой функции.

Если reduce передан список из одного элемента, то reduce вернёт этот элемент вне зависимости от переданной функции, например:
(reduce #'reduce '(1)) ; 1
\end{lstlisting}

\begin{lstlisting}
; 9. * Пусть list-of-list список, состоящий из списков. Написать функцию, которая вычисляет сумму длин всех элементов list-of-list (количество атомов), т.е. например для аргумента ((1 2) (3 4)) -> 4.
(defun atom-count-helper (lst cnt)
  (cond ((null lst) cnt)
        ((atom (car lst)) (atom-count-helper (cdr lst) (+ 1 cnt)))
        (t (atom-count-helper (cdr lst) cnt))))
(defun ac (lst)
  (atom-count-helper lst 0))

(defun f9-helper (lol cnt)
  (cond ((null lol) cnt)
        (t (f9-helper (cdr lol) (+ cnt (ac (car lol)))))))

(defun f9 (lol)
  (f9-helper lol 0))
\end{lstlisting}

% -------------------------------------------------------------------

\end{document}
